%% kapa: exam=1 round=23 year=2012 period=2 form=A subject=appraisal_law subject_ko=부동산관계법규 seq=38 number=78
%% source: Q-Net 감정평가사 기출문제 (cst00309) artlSeq=5052580 file=raw/1008/5052580/2012년 제23회 감정평가사 제1차시험 기출문제(2교시 A형).pdf
%% conversion: kordoc 4.24.51 -> tools/gen1.py (auto); review=proofread-sonnet; flags=none
\begin{question}{78}
부동산등기법상 건물의 표시에 관한 등기에 대한 설명으로 옳지 않은 것은?
\begin{choices}
  \item 건물이 멸실된 경우 그 건물 소유권의 등기명의인은 그 사실이 있는 때부터 1개월 이내에 멸실등기를 신청하여야 한다.
  \item 존재하지 아니하는 건물에 대한 등기가 있을 때에는 그 소유권의 등기명의인은 지체 없이 멸실등기를 신청하여야 한다.
  \item 건물의 합병이 있는 경우에 그 건물 소유권의 등기명의인은 그 사실이 있는 때부터 1개월 이내에 변경등기를 신청하여야 한다.
  \item 구분건물의 등기기록 중 1동 표제부에 관한 변경등기는 그 구분건물과 같은 1동의 건물에 속하는 다른 구분건물에 대하여는 변경등기로서의 효력이 없다.
  \item 구분건물로서 그 건물이 속하는 1동 전부가 멸실된 경우에 그 구분건물의 소유권의 등기명의인은 1동의 건물에 속하는 다른 구분건물의 소유권의 등기명의인을 대위하여 1동 전부에 대한 멸실등기를 신청할 수 있다.
\end{choices}
\end{question}
\answer{4}
